\uuid{cOi9}
\titre{Bornes larges ou strictes}
\auteur{Maxime Nguyen}
\organisation{AMSCC}
\datecreate{2026-11-24}
\niveau{L2}
\module{Probabilité}
\chapitre{Probabilité continue}
\sousChapitre{Densité de probabilité}
\difficulte{1}
\theme{densité, fonction de répartition, probabilité d'un intervalle}

\contenu{
\begin{question}
Soit $X$ une variable à densité $f$, de fonction de répartition $F$, et $a < b$. Quelles égalités sont vraies ?
\begin{answers}
\good{$\mathbb{P}(a \leq X \leq b) = \mathbb{P}(a < X < b)$}
\good{$\mathbb{P}(a \leq X \leq b) = F(b) - F(a)$}
\bad{$\mathbb{P}(a \leq X \leq b) = f(b) - f(a)$}
\good{$\mathbb{P}(a \leq X \leq b) = \displaystyle\int_a^b f(x)\,\mathrm{d}x$}
\bad{$\mathbb{P}(a \leq X \leq b) = F(b) - F(a) + f(a)$}
\end{answers}
\begin{explanations}
A est vrai parce que $\mathbb{P}(X=a) = \mathbb{P}(X=b) = 0$ ; ce n'est pas
vrai pour une variable discrète.

E est faux : la formule générale est $F(b) - F(a) + \mathbb{P}(X = a)$, et ici
$\mathbb{P}(X = a) = 0$. Remplacer $\mathbb{P}(X = a)$ par $f(a)$ revient à
prendre la densité pour une probabilité.

Contre-sens visé : transposer les réflexes du discret (« inclure la borne
ajoute de la probabilité ») ; confondre $F$ et $f$ dans un calcul de probabilité.
\end{explanations}
\end{question}
}
